\documentclass[11pt]{article}
\usepackage[T1]{fontenc}
\usepackage[utf8]{inputenc}
\usepackage{amsmath,amssymb,amsthm}
\usepackage[a4paper,margin=1in]{geometry}
\usepackage[hidelinks]{hyperref}
\setlength{\emergencystretch}{3em}

\newtheorem{theorem}{Theorem}
\newtheorem{lemma}[theorem]{Lemma}
\newcommand{\choosef}[2]{\binom{#1}{#2}}

\title{Erd\H{o}s Problem 700(ii):\\
An Unconditional Infinite Family}
\author{Nanocodex Erd\H{o}s 700 project}
\date{3 August 2026}

\begin{document}
\maketitle

\begin{abstract}
For
\[
 f(n)=\min_{2\leq k\leq\lfloor n/2\rfloor}
       \gcd\!\left(n,\choosef nk\right),
\]
we prove that infinitely many composite integers $n$ satisfy
$f(n)^2>n$.  The proof has two parts.  A Lucas-theorem argument shows that
$f(pqr)=pq$ for suitable nearby primes $p<q<r$.  The prime number theorem,
combined with an elementary packing lemma, supplies unboundedly many such
triples.  This exact result is kernel-checked in Lean.  This document makes
no novelty or priority claim; the purpose is to expose the repository's proof
in a compact, independently reviewable form.
\end{abstract}

\section{Lucas's criterion}

For a prime $s$, write $u\preceq_s v$ if every base-$s$ digit of $u$ is at
most the corresponding digit of $v$.  Lucas's theorem gives
\[
 s\nmid\choosef nk\quad\Longleftrightarrow\quad k\preceq_s n.
\]
When $s\mid n$, the units digit of $n$ is zero, and therefore
\[
 s\nmid\choosef nk
 \quad\Longleftrightarrow\quad
 s\mid k\ \text{ and }\ k/s\preceq_s n/s.                 \tag{1}
\]

\section{The nearby-prime structural lemma}

\begin{lemma}\label{lem:structural}
Let
\[
 q=p+a,\qquad r=q+c=p+a+c,\qquad b=a+c,
\]
where $p<q<r$ are primes, $0<a<c$, and $p>4b^3$.  Put $n=pqr$.
For every $0<k\leq n/2$, at most one of $p,q,r$ fails to divide
$\choosef nk$.
\end{lemma}

\begin{proof}
The hypotheses make the following genuine base expansions, with the units
digit written first:
\[
 qr=ab+(a+b)p+p^2,                                         \tag{2}
\]
\[
 pr=(q-ac)+(c-a-1)q+q^2,                                  \tag{3}
\]
\[
 pq=bc+(p-c)r.                                             \tag{4}
\]
All displayed digits are nonnegative and smaller than their bases because
$a,c<b$ and $p>4b^3$.

Suppose first that both $p$ and $q$ are absent from $\choosef nk$.  By (1),
$k=pqt$ for some positive integer $t$, and
\[
 qt\preceq_p qr,\qquad pt\preceq_q pr.                    \tag{5}
\]
The half-range gives $t\leq r/2$.  Write $at=up+v$ with $0\leq v<p$.
Then
\[
 qt=v+(t+u)p,
\]
and $t+u<p$: indeed $t<p$, $u<a$, and
$t+u\leq(p+b)/2+a-1<p$.  The next-to-units comparison with (2) gives
$t+u\leq a+b$.  Hence $at<p$; the units comparison then gives
$at\leq ab$, so $t\leq b$.  Thus $at<q$ and
\[
 pt=(q-at)+(t-1)q.
\]
Comparison with (3) gives $t\geq c$ from the units digit and
$t\leq c-a$ from the next digit, a contradiction.

Suppose next that $q$ and $r$ are both absent.  Write $k=qrt$.  Then
\[
 rt\preceq_q pr,\qquad qt\preceq_r pq,
\]
and $t\leq p/2$.  Write $ct=uq+v$, $0\leq v<q$.  Since $t+u<q$, the first
two base-$q$ digits of $rt$ are
\[
 rt=v+(t+u)q.
\]
Comparison with (3) gives $t+u\leq c-a-1$, and hence
$t\leq c-a-1$.  In particular $ct<r$, so
\[
 qt=(r-ct)+(t-1)r.
\]
The units comparison with (4) would give $r-ct\leq bc$.  But
$ct+bc<2b^2<p<r$, a contradiction.

Finally suppose that $p$ and $r$ are both absent.  Write $k=prt$.  Then
\[
 rt\preceq_p qr,\qquad pt\preceq_r pq,
\]
and $t\leq q/2$.  Write $bt=up+v$, $0\leq v<p$.  Again $t+u<p$, so the
first two base-$p$ digits of $rt$ are $v,t+u$.  Comparison with (2) gives
first $t\leq a+b$, then $bt<p$, then $bt\leq ab$, and therefore $t\leq a$.
Now
\[
 pt=(r-bt)+(t-1)r.
\]
Comparison with (4) would imply $r-bt\leq bc$, hence
$r\leq ab+bc=b^2$, contrary to $r>p>4b^3$.  All three possible pairs lead
to contradictions.
\end{proof}

\begin{lemma}\label{lem:value}
Under the hypotheses of Lemma~\ref{lem:structural},
\[
 f(pqr)=pq\qquad\text{and}\qquad f(pqr)^2>pqr.
\]
\end{lemma}

\begin{proof}
The integer $n=pqr$ is squarefree.  Lemma~\ref{lem:structural} says that
every admissible gcd contains at least two of $p,q,r$, and is therefore at
least $pq$.  At $k=r$, (4) gives $pq\bmod r=bc\geq1$, so
$1\preceq_r pq$ and criterion (1) shows that $r$ is absent.  The primes
$p$ and $q$ cannot be absent because either absence would force the distinct
prime to divide $k=r$.  Thus the gcd at $k=r$ is exactly $pq$, proving
$f(pqr)=pq$.

Finally
\[
 pq-r=p^2+(a-1)p-b>0
\]
under $p>4b^3$.  Hence $pq>r$ and
$(pq)^2>pqr$.
\end{proof}

\section{Producing infinitely many triples with the PNT}

\begin{lemma}[Asymmetric gaps in a short set]\label{lem:gaps}
If integers
$x_0<x_1<\cdots<x_{m-1}$ lie in an interval of length $T$ and no
$i<j<k$ satisfies
\[
 x_k-x_j>x_j-x_i,                                          \tag{6}
\]
then $2^{m-2}<T$.
\end{lemma}

\begin{proof}
Put $S_j=x_{m-1}-x_j$.  Failure of (6) for $(j,j+1,m-1)$ gives
$S_{j+1}\leq x_{j+1}-x_j$, and therefore
\[
 S_j=(x_{j+1}-x_j)+S_{j+1}\geq2S_{j+1}.
\]
Since $S_{m-2}\geq1$, iteration gives
$x_{m-1}-x_0=S_0\geq2^{m-2}$, while the interval hypothesis makes the span
strictly less than $T$.
\end{proof}

\begin{theorem}\label{thm:main}
There are infinitely many composite integers $n$ such that $f(n)^2>n$.
\end{theorem}

\begin{proof}
Let $T$ tend through sufficiently large positive integers and put
$N=8T^3$.  Partition $[N,2N)$ into $8T^2$ half-open intervals of length
$T$.  The prime number theorem gives
\[
 \pi(2N)-\pi(N)\sim\frac{N}{\log N}.
\]
Since
\[
 \frac{N/\log N}{8T^2(\log_2T+2)}\longrightarrow\infty,
\]
some length-$T$ interval contains more than
$\lfloor\log_2T\rfloor+2$ primes.  Lemma~\ref{lem:gaps} then supplies
primes $p<q<r$ in that interval with
\[
 r-q>q-p.
\]
Put $a=q-p$, $c=r-q$, and $b=a+c=r-p$.  Then $0<a<c$, $b<T$, and
\[
 4b^3<4T^3<8T^3\leq p.
\]
Lemma~\ref{lem:value} gives $f(pqr)^2>pqr$.  Moreover $pqr\geq(8T^3)^3$,
so these products are unbounded as $T\to\infty$.  An unbounded set of
natural numbers is infinite.
\end{proof}

\section{Formal verification and claim boundary}

The exact Lean declaration is
\path{Erdos700PNT.erdos_700_ii} in
\href{https://github.com/gakonst/nanocodex-erdos700/blob/main/lean-part2/Solution.lean}
{\texttt{lean-part2/Solution.lean}}, exposed through
\href{https://github.com/gakonst/nanocodex-erdos700/blob/main/lean-part2/Erdos700PNT.lean}
{\texttt{lean-part2/Erdos700PNT.lean}}.  The three pair-omission cases in
Lemma~\ref{lem:structural} are checked separately.  The repository verifier
rejects placeholders and reports only \texttt{propext},
\texttt{Classical.choice}, and \texttt{Quot.sound}, with no
\texttt{sorryAx} dependency.

This establishes the proposition represented in the repository.  It does
not by itself establish novelty, priority, publication, or community
acceptance.  The original problem is from P.~Erd\H{o}s and G.~Szekeres,
``Some number theoretic problems on binomial coefficients,''
\emph{Australian Mathematical Society Gazette} 5 (1978), 97--99,
\href{https://www.renyi.hu/~p_erdos/1978-46.pdf}{archival PDF}.

\end{document}
